\begin{blocksection}
\textbf{Data Modeling}

Data modeling is the process of deciding how your program will represent information. Ideally, you would have a single source of truth, meaning that storing each aspect of the state of your program in one place. If a list and the length of a list are both stored, which one is right? You can store what happened rather than derived consequences, which will make programs easier to test. In order to avoid invalid data, you can use a type system. Ideally, your types can only represent valid states. You can help ensure validity by using a small amount of functions to alter data. Finally, you can structure your data into types that functions will consume or return. Each function should take only the input it needs, not the whole state.

\vspace{\baselineskip}
\textbf{Web Pages}

A browser displays a page using three kinds of files. HTML describes what is on the page as a tree of elements, such as \lstinline$<div>$ and \lstinline$<button>$. CSS describes how each element is displayed: its size, color, font, and position. JavaScript is run by the browser and can change the contents of the page and respond to user actions. Gleam compiles to JavaScript.

A page is a tree of elements, and an element's children are written between its opening tag and closing tag. Nodes in this tree have attributes that control their style (via CSS). For example, an element with \lstinline$class="cell"$ is styled by the CSS rules for that class. Events are associated with nodes in this tree; for example, a click is an event on a particular button node. Changing the page involves changing this tree, and when the tree changes, the browser displays the new tree.

\vspace{\baselineskip}
\textbf{Model-View-Update (Lustre)}

A Lustre app has three parts. The model is the state of the app. The view is a function of the model that returns the tree of elements to display, and it specifies which message to send when the user takes an action. The update is a function that takes the current model and a message, and returns an updated model.

Lustre runs the app in a loop: it calls \lstinline$view$ on the model and displays the result. For every message that comes in, Lustre calls \lstinline$update$ with the current model and that message, then calls \lstinline$view$ on the new model and efficiently updates the page (by changing the tree) to match what \lstinline$view$ returns.

\vspace{\baselineskip}
\textbf{Building Elements}

Each function in \lstinline$lustre/element/html$, such as \lstinline$html.div$ and \lstinline$html.button$, creates a node in the tree. It takes a list of attributes and a list of children. For example, this expression is a \lstinline$div$ with one attribute (its class) and one child, a \lstinline$button$ with no attributes whose only child is the text ``Hi'':

\begin{lstlisting}
html.div([attribute.class("row")], [html.button([], [html.text("Hi")])])
\end{lstlisting}

It describes the HTML \lstinline$<div class="row"><button>Hi</button></div>$.

The function \lstinline$attribute.class("name")$ lets CSS style the node, and \lstinline$attribute.disabled(True)$ makes a button that can't be clicked. The function \lstinline$event.on_click(msg)$ associates a click event with the node, which sends \lstinline$msg$ to \lstinline$update$. A text child is created with \lstinline$html.text(s)$.

\lstinline$element.to_string$ converts an element to HTML text. Events are associated with nodes in the tree but are not part of the HTML text, so they don't appear. Layout can come from CSS rather than the shape of the tree.
\end{blocksection}